\documentclass[a4paper,11pt]{article}
\usepackage{kotex}
\usepackage{fontspec}
\usepackage{unicode-math}
\usepackage{amsmath,amssymb,amsthm}
\usepackage{array,booktabs,multirow,tabularx,colortbl}
\usepackage{enumitem}
\usepackage{geometry}
\usepackage{graphicx}
\usepackage{xcolor}
\usepackage{mdframed}
\usepackage{fancyhdr}
\usepackage{lastpage}
\usepackage{setspace}
\usepackage{tcolorbox}
\tcbuselibrary{skins,breakable}

\setmainfont{Apple SD Gothic Neo}
\setmathfont{Latin Modern Math}

\geometry{a4paper, top=15mm, bottom=20mm, left=15mm, right=15mm}

\setstretch{1.3}

\pagestyle{fancy}
\fancyhf{}
\renewcommand{\headrulewidth}{0pt}
\fancyfoot[C]{\small --- \thepage\ ---}

\newcounter{problemno}

\newcommand{\calcbox}[1][40mm]{%
  \vspace{2mm}
  \noindent\textbf{\small 풀이 과정}\\[-1mm]
  \noindent\rule{\linewidth}{0.3pt}\\[-2.5mm]
  \rule{\linewidth}{0.3pt}\\[-2.5mm]
  \rule{\linewidth}{0.3pt}\\[-2.5mm]
  \rule{\linewidth}{0.3pt}\\[-2.5mm]
  \rule{\linewidth}{0.3pt}
  \vspace{1mm}}

\newcommand{\essaybox}[1][5]{%
  \vspace{2mm}
  \noindent\rule{\linewidth}{0.3pt}\\[-2.5mm]
  \rule{\linewidth}{0.3pt}\\[-2.5mm]
  \rule{\linewidth}{0.3pt}\\[-2.5mm]
  \rule{\linewidth}{0.3pt}\\[-2.5mm]
  \rule{\linewidth}{0.3pt}
  \par\vspace{2mm}}

\definecolor{badgegray}{RGB}{60,60,60}
\definecolor{stepbg}{RGB}{245,245,245}
\definecolor{answerbg}{RGB}{235,248,235}

\begin{document}
\begin{center}
{\small\color{gray} 고2학년 1학기 인공지능수학}
\\[2mm]
{\LARGE\bfseries 텍스트 벡터화 연습}
\\[1mm]
{\normalsize\color{gray} TF · IDF · TF-IDF}
\end{center}
\vspace{-2mm}
\noindent\rule{\linewidth}{0.8pt}
\vspace{1mm}
\begin{tabular}{llll}
\textbf{학년} & \textbf{반} & \textbf{번} & \textbf{이름} \\
\end{tabular}
\par\vspace{2mm}
\noindent
\colorbox{badgegray}{\textcolor{white}{\small\bfseries 상황}}\enspace\textbf{영화 후기 3개를 수학적으로 분석해 봅시다.}
\vspace{1mm}
\begin{tcolorbox}[colback=white,colframe=black!70,boxrule=0.5pt,left=6pt,right=6pt,top=4pt,bottom=4pt]
아래 영화 후기에서 각 단어의 중요도를 TF-IDF로 계산하시오.\\[1mm]
\noindent 후기 A :  재미있는 영화. 영화 스토리가 정말 재미있고 좋아요.\\
\noindent 후기 B :  연기가 좋아요. 음악도 좋습니다.\\
\noindent 후기 C :  영화와 연기 모두 좋은 작품.\\
\end{tcolorbox}
\vspace{2mm}
\noindent
\colorbox{badgegray}{\textcolor{white}{\small\bfseries STEP 1}}\enspace\textbf{불용어 제거 + 단어 사전 U 만들기}
\par\vspace{1mm}
\begin{tcolorbox}[colback=stepbg,colframe=gray!50,boxrule=0.3pt,left=5pt,right=5pt,top=2pt,bottom=2pt]
\noindent\small 불용어 제거: 조사(가, 도, 와), 어미(요, 습니다, 는, 은, 이) 등 제거\\
\noindent\small 기본형 변환: 재미있는·재미있어요 → 재미있다  /  좋아요·좋습니다·좋은 → 좋다\\
\end{tcolorbox}
\vspace{1mm}
\noindent\textbf{후기 A 처리 결과 (예시)} : \enspace
\colorbox{stepbg}{\{ 재미있다, 영화, 스토리 \}}
\\[3pt]
\noindent\textbf{후기 B 처리 결과} : \enspace
\dotfill
\\[3pt]
\noindent\textbf{후기 C 처리 결과} : \enspace
\dotfill
\\[3pt]
\noindent\textbf{단어 사전 U} : \enspace
\dotfill
\quad\small\color{gray}(총 \hspace{25mm} 개, n = 3)
\\[3pt]
\par\vspace{3mm}
\noindent
\colorbox{badgegray}{\textcolor{white}{\small\bfseries STEP 2}}\enspace\textbf{TF 표 작성}\enspace{\small\color{gray} (각 후기에서 단어가 등장한 횟수)}
\par\vspace{1mm}
\begin{center}
\renewcommand{\arraystretch}{1.4}
\begin{tabular}{|c|c|c|c|c|c|c|c|}
\hline
\cellcolor{badgegray}\textcolor{white}{\small\bfseries 후기} & \cellcolor{badgegray}\textcolor{white}{\small\bfseries 영화} & \cellcolor{badgegray}\textcolor{white}{\small\bfseries 재미있다} & \cellcolor{badgegray}\textcolor{white}{\small\bfseries 스토리} & \cellcolor{badgegray}\textcolor{white}{\small\bfseries 연기} & \cellcolor{badgegray}\textcolor{white}{\small\bfseries 좋다} & \cellcolor{badgegray}\textcolor{white}{\small\bfseries 음악} & \cellcolor{badgegray}\textcolor{white}{\small\bfseries 작품} \\ \hline
\cellcolor{stepbg}\bfseries A & \cellcolor{stepbg} 2 & \cellcolor{stepbg} 2 & \cellcolor{stepbg} 1 &  &  &  &  \\ \hline
\cellcolor{stepbg}\bfseries B &  &  &  &  &  &  &  \\ \hline
\cellcolor{stepbg}\bfseries C &  &  &  &  &  &  &  \\ \hline
\end{tabular}
\end{center}
\par\vspace{3mm}
\noindent
\colorbox{badgegray}{\textcolor{white}{\small\bfseries STEP 3}}\enspace\textbf{DF와 IDF 계산}\enspace{\small\color{gray} (IDF = n ÷ DF,  n = 3)}
\par\vspace{1mm}
\begin{tcolorbox}[colback=stepbg,colframe=gray!50,boxrule=0.3pt,left=5pt,right=5pt,top=2pt,bottom=2pt]
\noindent\small DF: 해당 단어가 등장하는 후기의 수  |  IDF: 전체 후기 수 ÷ DF — 희귀한 단어일수록 IDF가 큽니다.\\
\end{tcolorbox}
\vspace{1mm}
\begin{center}
\renewcommand{\arraystretch}{1.4}
\begin{tabular}{|c|c|c|c|c|c|c|c|}
\hline
\cellcolor{badgegray}\textcolor{white}{\small\bfseries 단어} & \cellcolor{badgegray}\textcolor{white}{\small\bfseries 영화} & \cellcolor{badgegray}\textcolor{white}{\small\bfseries 재미있다} & \cellcolor{badgegray}\textcolor{white}{\small\bfseries 스토리} & \cellcolor{badgegray}\textcolor{white}{\small\bfseries 연기} & \cellcolor{badgegray}\textcolor{white}{\small\bfseries 좋다} & \cellcolor{badgegray}\textcolor{white}{\small\bfseries 음악} & \cellcolor{badgegray}\textcolor{white}{\small\bfseries 작품} \\ \hline
\cellcolor{stepbg}\bfseries DF &  &  &  &  &  &  &  \\ \hline
\cellcolor{stepbg}\bfseries IDF
(3÷DF) &  &  &  &  &  &  &  \\ \hline
\end{tabular}
\end{center}
\noindent\small $\Rightarrow$ "좋다"의 IDF = \hspace{25mm}     "재미있다"의 IDF = \hspace{25mm}     더 희귀한(특정 후기만의) 단어는? \hspace{25mm}
\par
\par\vspace{3mm}
\noindent
\colorbox{badgegray}{\textcolor{white}{\small\bfseries STEP 4}}\enspace\textbf{TF·IDF 계산 (후기 A)}\enspace{\small\color{gray} (= TF × IDF)}
\par\vspace{1mm}
\begin{center}
\renewcommand{\arraystretch}{1.4}
\begin{tabular}{|c|c|c|c|c|c|c|c|}
\hline
\cellcolor{badgegray}\textcolor{white}{\small\bfseries 단어} & \cellcolor{badgegray}\textcolor{white}{\small\bfseries 영화} & \cellcolor{badgegray}\textcolor{white}{\small\bfseries 재미있다} & \cellcolor{badgegray}\textcolor{white}{\small\bfseries 스토리} & \cellcolor{badgegray}\textcolor{white}{\small\bfseries 연기} & \cellcolor{badgegray}\textcolor{white}{\small\bfseries 좋다} & \cellcolor{badgegray}\textcolor{white}{\small\bfseries 음악} & \cellcolor{badgegray}\textcolor{white}{\small\bfseries 작품} \\ \hline
\cellcolor{stepbg}\bfseries TF(A) & \cellcolor{stepbg} 2 & \cellcolor{stepbg} 2 & \cellcolor{stepbg} 1 & \cellcolor{stepbg} 0 & \cellcolor{stepbg} 1 & \cellcolor{stepbg} 0 & \cellcolor{stepbg} 0 \\ \hline
\cellcolor{stepbg}\bfseries IDF & \cellcolor{stepbg}{\tiny\color{gray} ↑Step3} & \cellcolor{stepbg}{\tiny\color{gray} ↑} & \cellcolor{stepbg}{\tiny\color{gray} ↑} & \cellcolor{stepbg}{\tiny\color{gray} ↑} & \cellcolor{stepbg}{\tiny\color{gray} ↑} & \cellcolor{stepbg}{\tiny\color{gray} ↑} & \cellcolor{stepbg}{\tiny\color{gray} ↑} \\ \hline
\cellcolor{stepbg}\bfseries TF·IDF &  &  &  &  &  &  &  \\ \hline
\end{tabular}
\end{center}
\noindent\small $\Rightarrow$ 후기 A의 핵심어 (TF·IDF 최대): \hspace{25mm}   →   후기 A의 주제는 "\hspace{25mm}"이다.
\par
\par\vspace{3mm}
\noindent
\colorbox{badgegray}{\textcolor{white}{\small\bfseries STEP 5}}\enspace\textbf{생각 정리}
\par\vspace{1mm}
\noindent\textbf{Q1. "좋다"는 세 후기 모두에 등장한다. TF·IDF 관점에서 "좋다"가 특정 후기의 핵심어가 되기 어려운 이유를 설명하시오.}

\essaybox
\vspace{2mm}
\noindent\textbf{Q2. TF만 사용했을 때의 한계는 무엇이고, TF·IDF가 이를 어떻게 보완하는가?}

\essaybox
\vspace{2mm}
\par\vspace{3mm}
\begin{tcolorbox}[colback=white,colframe=black!30,boxrule=0.4pt,left=5pt,right=5pt,top=3pt,bottom=3pt,title={\small 확인 체크리스트},coltitle=black,fonttitle=\bfseries\small]
\noindent $\square$ \small 불용어를 제거하고 기본형으로 변환할 수 있다.\\
\noindent $\square$ \small 단어 사전 U를 만들고 TF 표를 완성할 수 있다.\\
\noindent $\square$ \small DF와 IDF를 계산할 수 있다.\\
\noindent $\square$ \small TF·IDF로 후기의 핵심어를 찾을 수 있다.\\
\end{tcolorbox}
\newpage
\begin{center}
{\Large\bfseries 정답 및 풀이}
\end{center}
\noindent\rule{\linewidth}{0.8pt}
\vspace{3mm}
\noindent
\colorbox{badgegray}{\textcolor{white}{\small\bfseries STEP 1}}\enspace\textbf{불용어 제거 + 단어 사전 U 만들기}
\par\vspace{1mm}
\begin{tcolorbox}[colback=stepbg,colframe=gray!50,boxrule=0.3pt,left=5pt,right=5pt,top=2pt,bottom=2pt]
\noindent\small 불용어 제거: 조사(가, 도, 와), 어미(요, 습니다, 는, 은, 이) 등 제거\\
\noindent\small 기본형 변환: 재미있는·재미있어요 → 재미있다  /  좋아요·좋습니다·좋은 → 좋다\\
\end{tcolorbox}
\vspace{1mm}
\noindent\textbf{후기 A 처리 결과 (예시)} : \enspace
\colorbox{stepbg}{\{ 재미있다, 영화, 스토리 \}}
\\[3pt]
\noindent\textbf{후기 B 처리 결과} : \enspace
\colorbox{answerbg}{\{ 연기, 좋다, 음악 \}}
\\[3pt]
\noindent\textbf{후기 C 처리 결과} : \enspace
\colorbox{answerbg}{\{ 영화, 연기, 좋다, 작품 \}}
\\[3pt]
\noindent\textbf{단어 사전 U} : \enspace
\colorbox{answerbg}{\{ 영화, 재미있다, 스토리, 연기, 좋다, 음악, 작품 \}  →  총 7개}
\quad\small\color{gray}(총 \hspace{25mm} 개, n = 3)
\\[3pt]
\par\vspace{3mm}
\noindent
\colorbox{badgegray}{\textcolor{white}{\small\bfseries STEP 2}}\enspace\textbf{TF 표 작성}\enspace{\small\color{gray} (각 후기에서 단어가 등장한 횟수)}
\par\vspace{1mm}
\begin{center}
\renewcommand{\arraystretch}{1.4}
\begin{tabular}{|c|c|c|c|c|c|c|c|}
\hline
\cellcolor{badgegray}\textcolor{white}{\small\bfseries 후기} & \cellcolor{badgegray}\textcolor{white}{\small\bfseries 영화} & \cellcolor{badgegray}\textcolor{white}{\small\bfseries 재미있다} & \cellcolor{badgegray}\textcolor{white}{\small\bfseries 스토리} & \cellcolor{badgegray}\textcolor{white}{\small\bfseries 연기} & \cellcolor{badgegray}\textcolor{white}{\small\bfseries 좋다} & \cellcolor{badgegray}\textcolor{white}{\small\bfseries 음악} & \cellcolor{badgegray}\textcolor{white}{\small\bfseries 작품} \\ \hline
\cellcolor{stepbg}\bfseries A & \cellcolor{stepbg} 2 & \cellcolor{stepbg} 2 & \cellcolor{stepbg} 1 & \cellcolor{answerbg} 0 & \cellcolor{answerbg} 1 & \cellcolor{answerbg} 0 & \cellcolor{answerbg} 0 \\ \hline
\cellcolor{stepbg}\bfseries B & \cellcolor{answerbg} 0 & \cellcolor{answerbg} 0 & \cellcolor{answerbg} 0 & \cellcolor{answerbg} 1 & \cellcolor{answerbg} 2 & \cellcolor{answerbg} 1 & \cellcolor{answerbg} 0 \\ \hline
\cellcolor{stepbg}\bfseries C & \cellcolor{answerbg} 1 & \cellcolor{answerbg} 0 & \cellcolor{answerbg} 0 & \cellcolor{answerbg} 1 & \cellcolor{answerbg} 1 & \cellcolor{answerbg} 0 & \cellcolor{answerbg} 1 \\ \hline
\end{tabular}
\end{center}
\par\vspace{3mm}
\noindent
\colorbox{badgegray}{\textcolor{white}{\small\bfseries STEP 3}}\enspace\textbf{DF와 IDF 계산}\enspace{\small\color{gray} (IDF = n ÷ DF,  n = 3)}
\par\vspace{1mm}
\begin{tcolorbox}[colback=stepbg,colframe=gray!50,boxrule=0.3pt,left=5pt,right=5pt,top=2pt,bottom=2pt]
\noindent\small DF: 해당 단어가 등장하는 후기의 수  |  IDF: 전체 후기 수 ÷ DF — 희귀한 단어일수록 IDF가 큽니다.\\
\end{tcolorbox}
\vspace{1mm}
\begin{center}
\renewcommand{\arraystretch}{1.4}
\begin{tabular}{|c|c|c|c|c|c|c|c|}
\hline
\cellcolor{badgegray}\textcolor{white}{\small\bfseries 단어} & \cellcolor{badgegray}\textcolor{white}{\small\bfseries 영화} & \cellcolor{badgegray}\textcolor{white}{\small\bfseries 재미있다} & \cellcolor{badgegray}\textcolor{white}{\small\bfseries 스토리} & \cellcolor{badgegray}\textcolor{white}{\small\bfseries 연기} & \cellcolor{badgegray}\textcolor{white}{\small\bfseries 좋다} & \cellcolor{badgegray}\textcolor{white}{\small\bfseries 음악} & \cellcolor{badgegray}\textcolor{white}{\small\bfseries 작품} \\ \hline
\cellcolor{stepbg}\bfseries DF & \cellcolor{answerbg} 3 & \cellcolor{answerbg} 1 & \cellcolor{answerbg} 1 & \cellcolor{answerbg} 2 & \cellcolor{answerbg} 3 & \cellcolor{answerbg} 1 & \cellcolor{answerbg} 1 \\ \hline
\cellcolor{stepbg}\bfseries IDF
(3÷DF) & \cellcolor{answerbg} 1 & \cellcolor{answerbg} 3 & \cellcolor{answerbg} 3 & \cellcolor{answerbg} 1.5 & \cellcolor{answerbg} 1 & \cellcolor{answerbg} 3 & \cellcolor{answerbg} 3 \\ \hline
\end{tabular}
\end{center}
\noindent\small $\Rightarrow$ "좋다"의 IDF = \hspace{25mm}     "재미있다"의 IDF = \hspace{25mm}     더 희귀한(특정 후기만의) 단어는? \hspace{25mm}

\noindent\small\color{gray} "좋다" IDF = 1  /  "재미있다" IDF = 3  /  더 희귀한 단어: 재미있다 (스토리, 음악, 작품도 IDF=3)
\par\vspace{3mm}
\noindent
\colorbox{badgegray}{\textcolor{white}{\small\bfseries STEP 4}}\enspace\textbf{TF·IDF 계산 (후기 A)}\enspace{\small\color{gray} (= TF × IDF)}
\par\vspace{1mm}
\begin{center}
\renewcommand{\arraystretch}{1.4}
\begin{tabular}{|c|c|c|c|c|c|c|c|}
\hline
\cellcolor{badgegray}\textcolor{white}{\small\bfseries 단어} & \cellcolor{badgegray}\textcolor{white}{\small\bfseries 영화} & \cellcolor{badgegray}\textcolor{white}{\small\bfseries 재미있다} & \cellcolor{badgegray}\textcolor{white}{\small\bfseries 스토리} & \cellcolor{badgegray}\textcolor{white}{\small\bfseries 연기} & \cellcolor{badgegray}\textcolor{white}{\small\bfseries 좋다} & \cellcolor{badgegray}\textcolor{white}{\small\bfseries 음악} & \cellcolor{badgegray}\textcolor{white}{\small\bfseries 작품} \\ \hline
\cellcolor{stepbg}\bfseries TF(A) & \cellcolor{stepbg} 2 & \cellcolor{stepbg} 2 & \cellcolor{stepbg} 1 & \cellcolor{stepbg} 0 & \cellcolor{stepbg} 1 & \cellcolor{stepbg} 0 & \cellcolor{stepbg} 0 \\ \hline
\cellcolor{stepbg}\bfseries IDF & \cellcolor{stepbg}{\tiny\color{gray} ↑Step3} & \cellcolor{stepbg}{\tiny\color{gray} ↑} & \cellcolor{stepbg}{\tiny\color{gray} ↑} & \cellcolor{stepbg}{\tiny\color{gray} ↑} & \cellcolor{stepbg}{\tiny\color{gray} ↑} & \cellcolor{stepbg}{\tiny\color{gray} ↑} & \cellcolor{stepbg}{\tiny\color{gray} ↑} \\ \hline
\cellcolor{stepbg}\bfseries TF·IDF & \cellcolor{answerbg} 2 & \cellcolor{answerbg} 6 & \cellcolor{answerbg} 3 & \cellcolor{answerbg} 0 & \cellcolor{answerbg} 1 & \cellcolor{answerbg} 0 & \cellcolor{answerbg} 0 \\ \hline
\end{tabular}
\end{center}
\noindent\small $\Rightarrow$ 후기 A의 핵심어 (TF·IDF 최대): \hspace{25mm}   →   후기 A의 주제는 "\hspace{25mm}"이다.

\noindent\small\color{gray} 핵심어: 재미있다 (TF·IDF = 6으로 최대)  →  주제: "재미" 또는 "재미있는 영화"
\par\vspace{3mm}
\noindent
\colorbox{badgegray}{\textcolor{white}{\small\bfseries STEP 5}}\enspace\textbf{생각 정리}
\par\vspace{1mm}
\noindent\textbf{Q1. "좋다"는 세 후기 모두에 등장한다. TF·IDF 관점에서 "좋다"가 특정 후기의 핵심어가 되기 어려운 이유를 설명하시오.}

\noindent\small\color{gray} "좋다"의 DF=3 (전체 후기 모두 등장)이므로 IDF = 3÷3 = 1로 낮다. TF·IDF = TF × 1 = TF에 불과해 가중치 효과가 없다. 즉, 모든 문서에 공통으로 나타나는 단어는 특정 문서를 구별하는 핵심어가 될 수 없다.
\vspace{2mm}
\noindent\textbf{Q2. TF만 사용했을 때의 한계는 무엇이고, TF·IDF가 이를 어떻게 보완하는가?}

\noindent\small\color{gray} TF만 사용하면 "영화", "좋다" 같은 범용 단어가 높은 점수를 받아 문서의 고유 특징을 반영하지 못한다. TF-IDF는 IDF를 곱해 여러 문서에 공통으로 등장하는 단어의 가중치를 낮추고, 특정 문서에만 나타나는 단어의 가중치를 높여 각 문서의 핵심어를 구별한다.
\vspace{2mm}
\par\vspace{3mm}
\end{document}